\documentclass[10pt,a4paper]{amsart}
\usepackage[margin=19mm]{geometry}
\usepackage{amsmath,amssymb,mathtools}
\usepackage[scr=rsfs,cal=euler]{mathalfa}
\usepackage{xfrac}
\usepackage[unicode,hidelinks,bookmarks=false]{hyperref}
\newcommand{\ie}{i.e.\ }
\newcommand{\cf}{cf.\ }
\newcommand{\resp}{respectively\ }
\newcommand{\Subsection}{\subsection}
\providecommand{\nobreakdash}{\nobreak\hbox{-}}
\setlength{\emergencystretch}{2em}
\multlinegap0pt
\usepackage{bm}
\usepackage{bbm}
\usepackage{dsfont}
\usepackage{xspace}
\usepackage{cancel}
\usepackage{mathtools}
\usepackage{amsthm}
\makeatletter
\@ifundefined{lemma}{\newtheorem{lemma}{Lemma}}{}
\makeatother
\makeatletter
\expandafter\ifx\csname Proof\endcsname\relax
\newenvironment{Proof}[1][Proof]{\begin{proof}[#1]}{\end{proof}}	% for redefining proofs
\fi
\newcommand{\R}{\mathbb{R}}	% for reals
\def\label@in@display@noarg#1{\cref@old@label@in@display{#1}}	% remove braces  
\expandafter\ifx\csname proofpart\endcsname\relax
\newenvironment{proofpart}[1]
{\vspace{3pt}
\refstepcounter{proofpart}%
\par\textit{Part~\arabic{proofpart}:~#1}.\,}
{\smallskip}
\fi
\makeatother
\title[Heavy-ball dynamics with Hessian-driven damping for non-conv]{Heavy-ball dynamics with Hessian-driven damping for non-convex optimization under the {\L}ojasiewicz condition}
\author{Vassilis Apidopoulos, Vasiliki Mavrogeorgou, Theodoros G. Tsironis}
\date{Source: arXiv:2506.11705v1 (CC BY 4.0)}
\begin{document}
\maketitle

\noindent\emph{Source and licence.} Heavy-ball dynamics with Hessian-driven damping for non-convex optimization under the {\L}ojasiewicz condition. Version arXiv:2506.11705v1 (CC BY 4.0), distributed under the Creative Commons Attribution 4.0 International licence, as recorded in the arXiv licence metadata for this exact version and shown on the abstract page https://arxiv.org/abs/2506.11705v1. Full rights notice: licenses/arxiv-2506.11705-CC-BY-4.0.txt.\par\smallskip

\noindent\textit{Editorial scope.} Counted unit: the Lemma whose source label is \texttt{lemma nonlinear gronwall power p} in the article (source0.tex lines 2677--2686, source label \texttt{lemma nonlinear gronwall power p}), delivered together with the article's own complete original proof of it (source0.tex lines 2687--2724, one \texttt{proof} environment of 38 lines). No mathematics is written, repaired or re-derived: statement and proof are the article's own text line for line. Required context retained uncounted: none. Every definition, display and label of those lines is retained unchanged. Omitted from this package and not redistributed: 1--2676, 2725--3386. Nothing inside a retained excerpt is omitted or altered: every retained body is delivered exactly as the article's lines are written, so no omission receipt of any kind accompanies this package. Citations: the retained excerpts carry no citation command at all, so no bibliography entry of the article is transcribed and no reference is redistributed. Reference adaptation: none was needed and none was made; every reference inside the delivered unit resolves inside it, so no unresolved reference or citation is produced. Printed slips kept literal: no printed word, symbol, hypothesis or number of the counted statement or of its proof was altered. Layout and class: the article's own class is \texttt{amsart} with packages including fix-cm, geometry, amsmath, amssymb, amsfonts, amsthm, amsaddr, mathtools, mathalfa, dsfont, libertine, AlegreyaSans, fontenc, acronym; the wrapper is the lane's pinned \texttt{amsart} editorial wrapper at 10pt on A4, which restates the theorem-like environments the retained text needs and adds the article's own macro definitions that the retained text needs (\texttt{\textbackslash R}, \texttt{\textbackslash label}), with extra wrapper packages (bm, bbm, dsfont, xspace, cancel, mathtools). The delivered numbering is the unit's own. Assets: no retained span contains a figure environment, a graphics-inclusion command, a PDF-page inclusion, a table or a TikZ picture; no image file is copied into this package, the package declares no asset dependency and no image file is redistributed with it. Licence: version arXiv:2506.11705v1 of this article is distributed under the Creative Commons Attribution 4.0 International licence, as recorded in the arXiv licence metadata for this exact version and shown on the abstract page https://arxiv.org/abs/2506.11705v1. A full rights notice is delivered at licenses/arxiv-2506.11705-CC-BY-4.0.txt. Characters: every retained excerpt is pure ASCII, so the delivered engine, XeLaTeX with its default Latin Modern text font, reports no missing character.
\begin{lemma}\label{lemma nonlinear gronwall power p}
Let $p>1$, $C>0$ and $t_{0}>0$ be three constants, $I=[t_{0},+\infty)$ and $u$,$G$ $:I\longrightarrow \R$, with $G\in L^{1}_{loc}(I;\R_{+})$ and $u\in\mathcal{C}^{1}(I)$, such that for all $t\in I$:
	\begin{equation}\label{relation nonlinear gronwall appendix}
		\dot{u}(t) \leq G(t) - Cu^{p}(t)
	\end{equation}
	Then, for all $t\geq t_{0}$, it holds:
\begin{equation}\label{bound nonlinear gronwall appendix}
	u(t) \leq \left(\max\left\{ \left(\frac{1}{C(p-1)}\right)^{\frac{1}{p-1}}, t_{0}^{\frac{1}{p-1}}u(t_{0}) \right\} + \int_{t_{0}}^{t}r^{\frac{1}{p-1}}G(r)dr \right)t^{-\frac{1}{p-1}}
\end{equation}
\end{lemma}

\begin{proof} 
We start by defining the following sets: 
$I_{1}:=\left\{t\in I ~:~ u^{p-1}(t) \leq \frac{1}{C(p-1)t} \right\}\cup\{t_{0}\}$ and $I_{2}:= \left\{t\in I ~:~ u^{p-1}(t) > \frac{1}{C(p-1)t} \right\}$, so that $I=I_{1}\cup I_{2}$.

Firstly, note that for all $t\in I_{1}$, \eqref{bound nonlinear gronwall appendix} holds immediately true.
On the other hand, for all $t\in I_{2}$, it holds  $Cu^{p}(t) \geq \frac{1}{(p-1)t}u(t)$. Therefore, by multiplying \eqref{relation nonlinear gronwall appendix} with $t^{\frac{1}{p-1}}$, it follows
% \begin{equation}\label{nonlineargronwall1}
% 	\dot{u}(t) \leq G(t) - \frac{1}{(p-1)t}u(t)
% \end{equation}
% By multiplying \eqref{nonlineargronwall1} with $t^{\frac{1}{p-1}}$, it follows
\begin{equation}\label{eq proof nonlinear growall}
	t^{\frac{1}{p-1}}\dot{u}(t) +\frac{1}{p-1}t^{\frac{1}{p-1}-1}u(t) \leq t^{\frac{1}{p-1}}G(t)
\end{equation}

In addition, since $I_{2}$ is and open set in $I$ for all $t\in I_{2}$, there exists $\bar{t} \in I_{2}$ and $0<\varepsilon<t-t_{0}$, such that $t\in (\bar{t}-\varepsilon,\bar{t}+\varepsilon)\subset I_{2}$ (which without loss of generality is considered to be maximal, i.e. $\bar{t}-\varepsilon \in I_{1}$). 
By integrating \eqref{eq proof nonlinear growall} on $(\bar{t}-\varepsilon,t)\subset I_{2}$, it follows:
\begin{equation}
	\begin{aligned}
	t^{\frac{1}{p-1}}u(t) & \leq (\bar{t}-\varepsilon)^{\frac{1}{p-1}}u(\bar{t}-\varepsilon) + \int_{\bar{t}-\varepsilon}^{t}r^{\frac{1}{p-1}}G(s)ds\\
	& \leq 	\max\left\{t_{0}^{\frac{1}{p-1}}u(t_{0}), \left(C(p-1\right)^{-\frac{1}{p-1}}\right\} +\int_{t_{0}}^{t}s^{\frac{1}{p-1}}G(s) ds
	\end{aligned}
\end{equation}	
where the last inequality follows from the fact that $\bar{t} - \varepsilon \in I_{1}$ and that $G(t)\geq 0$.
%\off{
%We split the proof in two cases:
%\begin{itemize}
%\item If $u^{p-1}(t) \leq \frac{1}{C(p-1)t}$, then \eqref{bound nonlinear gronwall appendix} holds immediately true. 
%\item $u^{p-1}(t) \geq \frac{1}{C(p-1)t}$, then it holds $Cu^{p}(t) \geq \frac{p}{t}u(t)$. Therefore, from \eqref{relation nonlinear gronwall appendix}, it follows
%	\begin{equation}\label{nonlineargronwall2}
%		\dot{u}(t) \leq G(t) - \frac{1}{(p-1)t}u(t)
%	\end{equation}
%By multiplying \eqref{nonlineargronwall2} with $t^{\frac{1}{p-1}}$ and integrating on $[s,t]$, it follows:
%\begin{equation}
%	t^{\frac{1}{p-1}}u(t) \leq s^{\frac{1}{p-1}}u(s) + \int_{s}^{t}r^{\frac{1}{p-1}}G(r)dr
%\end{equation}	
%\end{itemize}	
%}	
\end{proof}

\end{document}
